\documentclass[10pt]{article}

% Page and typography
\usepackage[margin=0.8in]{geometry}
\usepackage{microtype}
\usepackage{lmodern}
\usepackage{parskip}
\setlength{\parindent}{0pt}

% Compact section spacing
\usepackage{titlesec}
\titlespacing*{\section}{0pt}{0.7ex}{0.5ex}
\titlespacing*{\subsection}{0pt}{0.5ex}{0.3ex}
\titleformat{\section}{\bfseries\normalsize}{}{0pt}{}
\titleformat{\subsection}{\bfseries\normalsize}{}{0pt}{}

% Lists and hyperlinks
\usepackage{enumitem}
\setlist{nosep,leftmargin=*,labelsep=0.5em}
\usepackage[hidelinks]{hyperref}
\usepackage{verbatim}

% Metadata macros
\newcommand{\coursename}{CSE 493S: Advanced Topics in Machine Learning}
\newcommand{\term}{Autumn 2025}
\newcommand{\duedate}{Tuesday, October 22, 2025 at 11:59 PM}

\usepackage[backend=biber, style=numeric]{biblatex}
\addbibresource{citations.bib} % Link to your .bib file


\begin{document}

{\large \coursename\ — Project Proposal}\
{\normalsize \term \hfill  \\ Due: \duedate}
\par\noindent\rule{\linewidth}{0.4pt}

\textbf{Project Title:} \underline{Comparison of Performance of Large Language Model Distillation Methodologies for Medical QA}

\section*{Team Information}
Shreyan Mitra, \texttt{s99s42m}, Undergrad
\\
Bryan Zhao, \texttt{bohanz04}, Undergrad
\\
Frederico Baldan, \texttt{fbaldan}, Grad
\\
Tim Avilov, \texttt{timchick}, Undergrad \\

\section*{Goals and Expected Outcomes}

We will compare four different knowledge distillation (KD) methods. These methods are:
\begin{enumerate}
    \item \textbf{Sequence-level SFT} -- Training the student directly on the teacher's text outputs
    \item \textbf{Logit KD} -- Using token-level KL divergence loss combined with cross-entropy
    \item \textbf{Chain-of-Thought (CoT) distillation} -- Teaching the student to generate reasoning steps like the teacher
    \item \textbf{Preference-based KD} -- Using DPO (Direct Preference Optimization) or KTO methods
\end{enumerate}

We will perform this comparison by constructing an unified dataset called \textbf{Med-DistillMix-120k} (containing 120,000 examples from medical sources). We will divide this dataset into training, validation, and holdout sets.We will then benchmark the performance of the resulting distilled models on industry-standard datasets (MedQA, MedMCQA, PubMedQA, and PubHealth), and compare the performances. These benchmarks will be completely separate from our training set to prevent data leakage.To ensure model distribution actually improved, we'll compute perplexity on a fixed, held‑out biomedical text corpus: a 10k‑passage sample of PubMed abstracts (MedPPL‑10k). It is separate from both the training pool and from the medical benchmarks. We create a new evaluation suite called \textbf{FidelityBench-Med}. The purpose of this suite would be to ensure teacher-student fidelity (did the student actually learn the teacher's behavior or just score similarly) and evidence-grounded faithfulness (to prevent hallucinations).Finally, ablation studies will be conducted to isolate the effect of each component or hyperparameter by varying one factor at a time under a fixed compute/data budget, ensuring fair, reproducible comparisons and robust, defensible conclusions.

In summary, at the end of the project, we expect to deliver the following: (1) A unified training dataset called \textbf{Med-DistillMix-120k} (2) Evaluation results on four medical benchmarks: MedQA, MedMCQA, PubMedQA, and PubHealth (3) Perplexity measurements on biomedical text to measure language modeling quality (4) A new evaluation suite called \textbf{FidelityBench-Med} that measures how faithfully models cite evidence and avoid hallucinations (5) Tables, plots, and ablation studies testing different settings (dataset size, temperature, alpha mixing ratios, LoRA rank) and (6) Code to run the experiment

Our primary research question is which of the four distillation methods produce the best results. Our secondary research questions are as follows:
What are the trade-offs between accuracy, faithfulness to evidence, and computational efficiency? When does Chain-of-Thought distillation (or any other distillation method for that matter) help versus hurt performance on medical questions? Does having lower KL divergence (being closer to the teacher at the token level) actually translate to better task performance?
\\
\section*{First Milestone (Due Nov 6)}

By the first milestone, we build the training set, implement two distillation methodologies, perform the training for these methodologies, benchmarking and time permitting, perform perplexity, ablation, and fidelity analysis.
In summary, we aim to submit the following:
(1) MedMCQA dev accuracy; MedQA dev accuracy (2) PubMedQA accuracy and macro-F1 (3) Perplexity on 10k PubMed abstracts.* (4) Fidelity: token-level KL vs.\ teacher, top-1 agreement.* (5) Total GPU-hours per run. (6) Ablations: effects of $T$ and $\alpha$.* (7) Training logs and result tables.
\\
The delivery of items suffixed with * is subject to time constraints

\textit{Implementation details described below are subject to change and may contain inaccuracies that have not yet been sorted.}

Our basic training plan for achieving our first milestone is described below:
\begin{itemize}
\item QLoRA: rank 16--32, learning rate $2\times10^{-4}$, 10k--15k steps — We attach small “adapter” pieces to the big model and only train those (the big model stays frozen and compressed to save memory). “Rank 16–32” is how large these adapters are, the learning rate is how quickly they change, and “10k–15k steps” is how many training batches we run; this makes fine-tuning feasible on a single 24–40GB GPU.

\item Logit KD sweeps: temperature $T \in {2,3,4}$, mix $\alpha \in {0.3,0.5,0.7}$ — We try a few easy settings to see what works best. “Temperature $T$” controls how strongly the teacher shares its preferences across choices (higher = softer, more informative hints), and “$\alpha$” controls how much we copy the teacher versus follow the usual training target; we pick the combo that scores best on validation.

\item Shared tokenizer; sequence length 1{,}024--1{,}536 tokens — We make both models split text into tokens the same way so inputs line up correctly, and we cap each example to about 1k–1.5k tokens so it fits in GPU memory while still keeping enough context for medical questions and explanations.
\end{itemize}

\section{Empirical + Theoretical Analysis}

While this is primarily an empirical comparison project, we'll add a theoretical-style analysis of fidelity -- specifically examining the mathematical relationship between KL divergence to the teacher and downstream performance.

We'll measure the correlation between token-level KL divergence / top-$k$ prediction overlap and both:
\begin{enumerate}
    \item Benchmark accuracy (how well the model answers questions)
    \item Faithfulness metrics (how well the model cites evidence)
\end{enumerate}

We'll fit simple regression models and report confidence intervals to understand whether being more faithful to the teacher actually matters for real-world performance.
\\
\newpage
\textbf{APPENDIX}
\\
\section*{Data and Datasets}
1. Public, non-PHI sources (licensed):
\begin{itemize}
\item MedQA-USMLE (medical exams): \url{https://huggingface.co/datasets/medqa/medqa-usmle}
\item MedMCQA (multi-subject MCQ): \url{https://huggingface.co/datasets/medmcqa/medmcqa}
\item PubMedQA (biomedical QA): \url{https://huggingface.co/datasets/pubmed_qa}
\item PubHealth (health claims): \url{https://huggingface.co/datasets/EmreCikara/pubhealth}
\end{itemize}

2. Training set: Med-DistillMix-120k (for full study)
\begin{itemize}
\item 80k MedMCQA/MedQA stems (prompts only), 20k PubMedQA, 20k PubHealth.
\item Cleaned/deduplicated; split 90% train / 5% val / 5% holdout; no overlap with official eval sets.
\item Artifacts: teacher responses (all); CoT rationales for 30k--40k reasoning prompts; DPO pairs (chosen = teacher, rejected = baseline/perturbed); optional cached top-$k$ logits (val/holdout) for speed.
\end{itemize}

3. Evaluation Benchmarks
\begin{itemize}
\item MedQA-USMLE, MedMCQA: accuracy.
\item PubMedQA, PubHealth: accuracy and macro-F1.
\item Language modeling: perplexity on 10k PubMed abstracts.
\item FidelityBench-Med (1k--2k prompts): evidence faithfulness (RAGAS/NLI), citation coverage, hallucination rate; teacher--student fidelity (token-level KL, top-$k$ overlap).
\end{itemize}

\section*{Models}
Teacher: Meditron-7B (primary): \url{https://huggingface.co/epfl-llm/meditron-7b}
\\
Student: Qwen2-1.5B (primary)

\section*{Code and Tools}
\begin{itemize}
\item Hugging Face Transformers/Datasets; TRL (SFTTrainer, DPOTrainer); PEFT (LoRA/QLoRA).
\item lm-eval or lighteval for benchmarks.
\item Simple BM25 retriever (Elasticsearch or PyTerrier) for evidence.
\end{itemize}

\section*{Compute and Feasibility}
\textbf{Hardware:} 1--2$\times$A100 40GB or 1$\times$RTX 4090 24GB.\
\\
\textbf{Config:} QLoRA rank 16--32, bf16; seq len 1{,}024--1{,}536; global batch 128 (via accumulation); 10k--20k steps/method.\
\\
\textbf{Estimated GPU-hours/method:}
\begin{itemize}
\item SFT: 25--40; Logit KD: 35--55 (teacher forward cost);
\item CoT: 30--50 (longer sequences); DPO: 30--45;
\item Total: $\sim$120--190 GPU-hours.
\end{itemize}

\newpage
\nocite{*}
\printbibliography

\end{document}

